% Part: methods
% Chapter: induction
% Section: structural-induction

\documentclass[../../../include/open-logic-section]{subfiles}

\begin{document}

\olfileid{mth}{ind}{sti}

\olsection{నిర్మాణాత్మక ఆగమనం}

ఇంతవరకు అన్ని సహజ సంఖ్యల గురించిన ఫలితాలను స్థాపించడానికి ఆగమనాన్ని వాడాం. అయితే, ఆగమనాత్మకంగా నిర్వచించిన సమితిలోని అన్ని \tetoken{మూలకాల}{!!{element}s} గురించిన ఫలితాలను నేరుగా నిరూపించడానికి దీనికి అనురూపమైన సూత్రాన్ని వాడవచ్చు. నిర్వచించిన వస్తువుల నిర్మాణంపై ఆధారపడుతుంది కనుక దీనిని తరచుగా \emph{నిర్మాణాత్మక} ఆగమనం అంటారు.

సాధారణంగా, ఆగమనాత్మక నిర్వచనం (a)~సమితిలోని ``ప్రారంభ'' \tetoken{మూలకాల}{!!{element}s} జాబితాను, (b)~పాత మూలకాల నుంచి కొత్త \tetoken{మూలకాలను}{!!{element}s} ఉత్పత్తి చేసే క్రియల జాబితాను ఇస్తుంది. ఉదాహరణకు, చక్కని పదాల విషయంలో ప్రారంభ వస్తువులు అక్షరాలు. ఇక్కడ ఒకే క్రియ ఉంది:
\begin{align*}
  o(s_1, s_2) = & [s_1 \circ s_2]
\end{align*}
సహజ సంఖ్యలు~$\Nat$ కూడా ఆగమనాత్మక నిర్వచనం ద్వారా వచ్చాయని భావించవచ్చు: ప్రారంభ వస్తువు~$0$; క్రియ అనువర్తి ప్రమేయం~$x + 1$.

ఆగమనాత్మకంగా నిర్వచించిన సమితిలోని అన్ని మూలకాల గురించి ఏదైనా నిరూపించాలంటే, అంటే సమితిలోని ప్రతి \tetoken{మూలకానికి}{!!{element}} $P$ అనే లక్షణం ఉందని నిరూపించాలంటే, మనం ఈ రెండు పనులు చేయాలి:
\begin{enumerate}
\item ప్రారంభ వస్తువులకు~$P$ ఉందని నిరూపించాలి.
\item ప్రతి క్రియ~$o$కు, దాని ఆర్గ్యుమెంట్లకు~$P$ ఉంటే దాని ఫలితానికీ ఉందని నిరూపించాలి.
\end{enumerate}
ఉదాహరణకు, అన్ని చక్కని పదాల గురించిన వాదనను నిరూపించాలంటే అది అన్ని అక్షరాలకూ సత్యమని, $s_1$, $s_2$లకు విడివిడిగా సత్యమైతే $[s_1 \circ s_2]$కూ సత్యమని చూపుతాం.

\begin{prop}
  ఏ చక్కని పదం~$t$లోనైనా $[$ల సంఖ్య $]$ల సంఖ్యకు సమానం.
\end{prop}

\begin{proof}
నిర్మాణాత్మక ఆగమనాన్ని వాడతాం. చక్కని పదాలు ఆగమనాత్మకంగా నిర్వచించబడ్డాయి: అక్షరాలు ప్రారంభ వస్తువులు; పాత పదాల నుంచి కొత్త చక్కని పదాలను నిర్మించడానికి $o$ అనే క్రియ ఉంది.
\begin{enumerate}
\item ప్రతి అక్షరానికీ వాదన సత్యం: ఒక్క అక్షరంలో $[$ల సంఖ్య~$0$, $]$ల సంఖ్య కూడా~$0$.
\item $s_1$లో $[$ల సంఖ్య $]$ల సంఖ్యకు సమానమని, $s_2$కూ అదే సత్యమని అనుకుందాం. $o(s_1, s_2)$లో, అంటే $[s_1 \circ s_2]$లో, $[$ల సంఖ్య $s_1$, $s_2$లలోని $[$ల సంఖ్యల మొత్తానికి ఒకటి కలిపినది. $o(s_1, s_2)$లో $]$ల సంఖ్య $s_1$, $s_2$లలోని $]$ల సంఖ్యల మొత్తానికి ఒకటి కలిపినది. అందువల్ల $o(s_1, s_2)$లో $[$ల సంఖ్య $o(s_1,s_2)$లో $]$ల సంఖ్యకు సమానం.
\end{enumerate}
\end{proof}

\begin{prob}
  ఏ చక్కని పదమూ~$]$తో మొదలుకాదని నిర్మాణాత్మక ఆగమనం ద్వారా నిరూపించండి.
\end{prob}

నిర్మాణాత్మక ఆగమనంతో మరో నిరూపణను చూద్దాం. చిహ్నాల శ్రేణి~$t$కు \emph{నిజ ప్రారంభ భాగం} అంటే, ఎడమ నుంచి చదివినప్పుడు చిహ్నం చిహ్నానికి $t$తో సరిపడుతూ, $t$ కన్నా చిన్నగా ఉండే శ్రేణి~$s$. ఉదాహరణకు, $[a \circ {}$ అనేది $[a \circ b]$కు నిజ ప్రారంభ భాగం; కానీ $[b \circ {}$ కాదు (రెండవ చిహ్నం వేరుగా ఉంది), $[a
\circ b]$ కూడా కాదు (అదే పొడవు).

\begin{prop}\ollabel{prop:initial}
  ఏ చక్కని పదం~$t$ యొక్క ప్రతి నిజ ప్రారంభ భాగంలో $]$ల కంటే
    $[$లే ఎక్కువ.
\end{prop}

\begin{proof}
  $t$పై ఆగమనం ద్వారా:
  \begin{enumerate}
  \item $t$ ఒక్క అక్షరం: అప్పుడు $t$కు నిజ ప్రారంభ భాగాలు లేవు.
  \item ఏవో చక్కని పదాలు $s_1$,~$s_2$ కోసం $t = [s_1 \circ s_2]$. $r$ అనేది $t$ యొక్క నిజ ప్రారంభ భాగమైతే ఈ సాధ్యాలు ఉన్నాయి:
    \begin{enumerate}
    \item $r$ కేవలం $[$: అప్పుడు $r$లో $]$ల కంటే ఒక $[$ ఎక్కువ.
    \item $r$ అనేది $[r_1$; ఇక్కడ $r_1$,~$s_1$ యొక్క నిజ ప్రారంభ భాగం: $s_1$ చక్కని పదం కనుక ఆగమన పరికల్పన ప్రకారం $r_1$లో $]$ల కంటే $[$లే ఎక్కువ; $[r_1$కు కూడా అదే నిజం.
    \item $r$ అనేది $[s_1$ లేదా $[s_1 \circ {}$: ముందరి ఫలితం ప్రకారం $s_1$లో $[$ల, $]$ల సంఖ్యలు సమానం; కాబట్టి $[s_1$లోనూ $[s_1 \circ {}$లోనూ $]$ల కంటే ఒక $[$ ఎక్కువ.
    \item $r$ అనేది $[s_1 \circ r_2$; ఇక్కడ $r_2$,~$s_2$ యొక్క నిజ ప్రారంభ భాగం: ఆగమన పరికల్పన ప్రకారం $r_2$లో $]$ల కంటే $[$లే ఎక్కువ. ముందరి ఫలితం ప్రకారం $s_1$లో $[$ల, $]$ల సంఖ్యలు సమానం. కాబట్టి $[s_1 \circ r_2$లో $]$ల కంటే $[$ల సంఖ్య ఎక్కువ.
    \item $r$ అనేది $[s_1 \circ s_2$: ముందరి ఫలితం ప్రకారం $s_1$లో $[$ల, $]$ల సంఖ్యలు సమానం; $s_2$కూ అదే నిజం. కనుక $[s_1 \circ s_2$లో $]$ల కంటే ఒక $[$ ఎక్కువ.
    \end{enumerate}
  \end{enumerate}
\end{proof}

\end{document}
